% !TEX program = lualatex
\documentclass[a4paper,10pt,twocolumn]{article}
\usepackage{kaitoushu}
\renewcommand{\baselinestretch}{1.2}
\lhead{応用数学 問題集}
\problemrange{問156〜160}

\begin{document}
\thispagestyle{fancy}

\begin{center}
  {\Large \textbf{応用数学　3章　第13回　Basic　解答}}
\end{center}
\vspace{0.5em}

\noindent\textbf{\large【Basic】}
\vspace{0.3em}

\mondai{156}{
次の公式を証明せよ．ただし，$a$ は $0$ でない実数とする．$\mathcal{F}[f(ax)] = \dfrac{1}{|a|}F\!\left(\dfrac{u}{a}\right)$
\hfill {\scriptsize [教p.97(問5)]}
}

\kotae{
$\mathcal{F}[f(ax)] = \displaystyle\int_{-\infty}^{\infty}f(ax)e^{-iux}\,dx$

$t=ax$ と置換（$dx=dt/a$）：

$a>0$：積分範囲変わらず，$\dfrac{1}{a}\displaystyle\int_{-\infty}^{\infty}f(t)e^{-i(u/a)t}\,dt = \dfrac{1}{a}F\!\left(\dfrac{u}{a}\right)$

$a<0$：積分範囲が逆転して符号変わり，$\dfrac{1}{a}\displaystyle\int_{+\infty}^{-\infty}f(t)e^{-i(u/a)t}\,dt = -\dfrac{1}{a}F\!\left(\dfrac{u}{a}\right)$

いずれも $\dfrac{1}{|a|}F\!\left(\dfrac{u}{a}\right)$
}

\mondai{157}{
問題153の $f(x)$ と $g(x)$ について，$\mathcal{F}[f(x) * g(x)]$ を求めよ．
\hfill {\scriptsize [教p.98(問6)]}
}

\kotae{
たたみこみ定理：$\mathcal{F}[f*g] = \mathcal{F}[f]\cdot\mathcal{F}[g]$

問153より $\mathcal{F}[f] = \dfrac{-2i(e^{3iu}-1)}{u}$，$\mathcal{F}[g] = \dfrac{(1+iu)e^{-iu}-1}{u^2}$

$\mathcal{F}[f*g] = \dfrac{-2i(e^{3iu}-1)}{u}\cdot\dfrac{(1+iu)e^{-iu}-1}{u^2} = \dfrac{-2i(e^{3iu}-1)\{(1+iu)e^{-iu}-1\}}{u^3}$
}

\mondai{158}{
公式 $\mathcal{F}[e^{-ax^2}] = \sqrt{\dfrac{\pi}{a}}\,e^{-\frac{u^2}{4a}}$ を用いて，次の関数のフーリエ変換を求めよ．
(1) $e^{-\frac{x^2}{4}}$ \quad (2) $xe^{-\frac{x^2}{4}}$ \quad (3) $x^2 e^{-\frac{x^2}{4}}$
\hfill {\scriptsize [教p.98(問7)]}
}

\kotae{
(1) $a=1/4$ として $\mathcal{F}[e^{-x^2/4}] = \sqrt{4\pi}\,e^{-u^2} = 2\sqrt{\pi}\,e^{-u^2}$

(2) 微分法則 $\mathcal{F}[xf(x)] = i\dfrac{d}{du}F(u)$ より
$\mathcal{F}[xe^{-x^2/4}] = i\cdot(-2u)\cdot 2\sqrt{\pi}\,e^{-u^2} = -4i\sqrt{\pi}\,ue^{-u^2}$

(3) $\mathcal{F}[x^2 e^{-x^2/4}] = (i)^2\dfrac{d^2}{du^2}(2\sqrt{\pi}\,e^{-u^2})$

$\dfrac{d^2}{du^2}(e^{-u^2}) = (-2+4u^2)e^{-u^2}$

$= -2\sqrt{\pi}(-2+4u^2)e^{-u^2} = 4\sqrt{\pi}(1-2u^2)e^{-u^2}$
}

\mondai{159}{
$\mathcal{F}^{-1}[e^{-3u^2}]$ を求めよ．
\hfill {\scriptsize [教p.98(問8)]}
}

\kotae{
$\mathcal{F}[e^{-ax^2}] = \sqrt{\dfrac{\pi}{a}}\,e^{-u^2/(4a)}$ で $a=3$，$u$ と $x$ の役割を入れ替えて

$\mathcal{F}^{-1}[e^{-3u^2}] = \dfrac{1}{2\sqrt{3\pi}}\,e^{-x^2/12}$
}

\mondai{160}{
次の等式を証明せよ．
(1) $\mathcal{F}\!\left[e^{-\frac{x^2}{4}} * xe^{-\frac{x^2}{4}}\right] = -8\pi iu\,e^{-2u^2}$
(2) $e^{-\frac{x^2}{4}} * xe^{-\frac{x^2}{4}} = \sqrt{\dfrac{\pi}{2}}\,xe^{-\frac{x^2}{8}}$
\hfill {\scriptsize [教p.99(問9)]}
}

\kotae{
(1) $\mathcal{F}[e^{-x^2/4}] = 2\sqrt{\pi}\,e^{-u^2}$，$\mathcal{F}[xe^{-x^2/4}] = -4i\sqrt{\pi}\,ue^{-u^2}$

$\mathcal{F}[e^{-x^2/4}*xe^{-x^2/4}] = 2\sqrt{\pi}\,e^{-u^2}\cdot(-4i\sqrt{\pi}\,ue^{-u^2}) = -8\pi iu\,e^{-2u^2}$

(2) $\mathcal{F}\!\left[\sqrt{\dfrac{\pi}{2}}\,xe^{-x^2/8}\right]$ を計算する．

$a=\dfrac{1}{8}$ として公式を適用：$\sqrt{\dfrac{\pi}{a}}=\sqrt{8\pi}=2\sqrt{2\pi}$，$e^{-u^2/(4\cdot 1/8)}=e^{-2u^2}$

$\mathcal{F}[e^{-x^2/8}] = 2\sqrt{2\pi}\,e^{-2u^2}$

微分法則 $\mathcal{F}[xf]=i\dfrac{d}{du}F(u)$：$\mathcal{F}[xe^{-x^2/8}] = i(-4u)\cdot 2\sqrt{2\pi}\,e^{-2u^2} = -8i\sqrt{2\pi}\,ue^{-2u^2}$

$\mathcal{F}\!\left[\sqrt{\dfrac{\pi}{2}}\,xe^{-x^2/8}\right] = \sqrt{\dfrac{\pi}{2}}\cdot(-8i\sqrt{2\pi}\,ue^{-2u^2}) = -8\pi iu\,e^{-2u^2}$

(1) の結果と一致するので，逆変換の一意性より等式が成立する．
}

\end{document}